\chapter{Stochastic framework and Uncertainty Quantification} \label{appendixB}
\textit{Note: The theoretical stochastic framework, reliability concepts, and numerical uncertainty quantification methods (including Gaussian Quadrature, Random Fields, and Sobol sensitivity) presented in this appendix are comprehensively synthesised from the foundational works of Baroth et al.~\cite{baroth2011, baroth2023} and Sudret et al.~\cite{sudret2000, sudret2003}.}

% =====================================================================
\section{Sources and classification of uncertainty} \label{sec:uncertainty-theory}

\subsection{Motivation for probabilistic assessment}

In civil engineering, a structure is considered safe when it remains fit for its intended use throughout its service life \cite{baroth2011}. Deterministic models have traditionally been used to evaluate structural safety. However, the gap between a real physical system and a simplified numerical model, along with the unforeseeable nature of future loads and environmental conditions, inherently leads to uncertainty \cite{baroth2023}.

To manage these risks, modern engineering codes such as the Eurocodes adopt a probabilistic format \cite{baroth2011}. This allows the safety of a structure to be explicitly quantified through its probability of failure $P_f$ (\cref{sec:pf-beta}) and its reliability $R = 1 - P_f$. By introducing random variables into mechanical models and propagating them through structural analyses (\cref{sec:stochastic-framework}), objective decisions can be made based on both the probability of an event and the severity of its consequences \cite{baroth2023}.

\subsection{Classification of uncertainties}

According to reliability theory, uncertainties can be classified into two categories \cite{baroth2011,baroth2023}:

\begin{itemize}
    \item \textbf{Aleatoric uncertainties} arise from the natural randomness of the environment and the spatial or temporal variability of material properties. They cannot be reduced; they can only be quantified and accounted for in the design process \cite{baroth2023}.

    \item \textbf{Epistemic uncertainties} stem from incomplete knowledge or insufficient data. They include \emph{model uncertainties} (discrepancy between mathematical models and actual physical processes) and \emph{statistical uncertainties} (limited or inaccurate data). Unlike aleatoric uncertainties, they can be reduced by improving models and gathering more data \cite{baroth2023}.
\end{itemize}

\subsection{Concrete heterogeneity and material variability}

A critical source of aleatoric uncertainty is the inherent heterogeneity of concrete \cite{baroth2023,delarrard2010}. Its mechanical properties depend on the formation history, casting conditions, internal chemical reactions, and long-term degradation processes \cite{delarrard2010}.

This variability manifests as \emph{continuous} spatial fluctuations (e.g.\ Young's modulus, permeability, porosity) and \emph{discrete} defects (e.g.\ micro-cracks, capillary networks) \cite{baroth2023}. Consequently, material parameters must be modelled as random variables or random fields rather than deterministic values (\cref{sec:stochastic-framework}).

\section{Random variables and distributions} \label{sec:random-variables}
In this stochastic framework, material properties are modelled as random variables characterized by their Probability Density Functions (PDF) and Cumulative Distribution Functions (CDF). The distributions are defined by their statistical moments (mean $\mu$, variance $\sigma^2$, skewness, and kurtosis) and the dimensionless Coefficient of Variation ($\mathrm{CoV} = \sigma/\mu$).

The concepts of PDF and CDF form the mathematical foundation for the uncertainty propagation methods used in this work: Latin Hypercube Sampling (LHS) relies on inverting the CDF for variance-reduced sampling, while Gaussian Quadrature analytically reconstructs the response PDF from its statistical moments to evaluate the probability of failure.

\textbf{Choice of probability distributions for concrete properties} \label{sec:distribution-justification}
In the stochastic assessment of concrete structures (as performed in \cref{chapter5}), the choice of probability distributions must reflect the physical nature of the uncertainty. Three main distribution families are typically employed:

\begin{description}
    \item[Lognormal --- permeability $K_0$.] The intrinsic permeability is strictly positive and varies by multiplicative factors: it is known only to within an order of magnitude and its scatter is naturally expressed as a ratio, not an additive offset. A lognormal law guarantees $K_0>0$ for every sample and reproduces this multiplicative, positively skewed variability; it is the standard model for the permeability of concrete \cite{delarrard2010,sudret2000}. It is also consistent with the random-field generator of \cref{sec:spatial-effect}: the Cast3M \texttt{ALEA} operator produces a Gaussian field that is exponentiated into a lognormal one (\cref{sec:lognormal-transformation}), so the random-variable and random-field descriptions of $K_0$ share the same marginal.
    
    \item[Normal --- cement content $c$ and aggregate content $agg$.] The batched masses (in \si{\kilogram\per\cubic\meter}) are controlled by the mixing plant around a target value. Their variability is the accumulation of many small, independent dosing and weighing errors, which by the central-limit theorem tends to a Gaussian centred on the nominal dosage. The scatter is small relative to the mean ($\mathrm{CoV}\approx 4\%$ for $c$, $\approx 3\%$ for $agg$), so the symmetric Normal law describes the batch-to-batch dispersion faithfully and the probability of a physically meaningless (negative) sample is negligible.
    
    \item[Uniform --- ratios $w/c$ and $s/c$.] For the water/cement and silica/cement ratios, what is known is not a measured central tendency but an \emph{admissible design window}: an HPC of this type is specified to lie within $w/c\in[0.30,0.38]$ and $s/c\in[0.08,0.12]$. When only the bounds are known and there is no basis to prefer any value inside them, the maximum-entropy (least-informative) choice is the uniform distribution \cite{baroth2023}. It models the \emph{epistemic} imprecision of the as-specified mix rather than a measurement scatter, and it keeps every sample strictly within the physically admissible window, avoiding tails that would push Powers' model or the solver outside its valid range.
\end{description}

\section{Random fields} \label{sec:random-fields}

A random variable assigns a \emph{single} uncertain value to a material property for the entire structure. In reality, concrete properties vary spatially due to heterogeneous microstructure, local variations in the water-to-cement ratio, and non-uniform compaction during casting \cite{delarrard2010}. A \textbf{random field} captures this spatial variability: it assigns an uncertain value to \emph{every point} $\mathbf{x}$ in the domain, not just one value for the whole structure.

Formally, a random field is a collection of random variables indexed by position \cite{sudret2000}:

\begin{equation}
    H \;=\; \bigl\{H(\mathbf{x},\omega) :
    \mathbf{x} \in \mathcal{D},\;
    \omega \in \Omega \bigr\},
    \label{eq:rf-definition}
\end{equation}

where $\omega$ denotes one particular realisation drawn from the probability space. For a fixed $\omega$, $H(\cdot,\omega)$ is a deterministic spatial function (one ``map'' of the property over the structure); for a fixed $\mathbf{x}$, $H(\mathbf{x},\cdot)$ is an ordinary random variable.

A \textbf{second-order} random field is fully characterised by two functions \cite{sudret2000}:

\begin{itemize}
    \item the \textbf{mean field}
    $\bar{H}(\mathbf{x}) = E\!\left[H(\mathbf{x})\right]$,
    \item the \textbf{covariance function}
    $C_H(\mathbf{x},\mathbf{y}) = E\!\left[
        \bigl(H(\mathbf{x}) - \bar{H}(\mathbf{x})\bigr)
        \bigl(H(\mathbf{y}) - \bar{H}(\mathbf{y})\bigr)
    \right]$.
\end{itemize}

The mean field gives the average value at each point; the covariance function measures how strongly the values at two points $\mathbf{x}$ and $\mathbf{y}$ are linked. When both depend only on the separation $\boldsymbol{\tau} = \mathbf{x} - \mathbf{y}$, the field is called \textbf{homogeneous} (stationary).

\subsection{Correlation function and correlation length} \label{sec:correlation}

When a material property (e.g.\ Young's modulus) is modelled as a random field, nearby points tend to have similar values while distant points are more different. The \textbf{autocorrelation function} $\rho_H(\boldsymbol{\tau})$ quantifies this similarity as a function of distance $\tau$:

\begin{equation}
    \rho_H(\boldsymbol{\tau})
    \;=\; \frac{C_H(\boldsymbol{\tau})}{\sigma_H^2},
    \qquad \rho_H(\mathbf{0}) = 1,
    \label{eq:autocorrelation}
\end{equation}

where $\sigma_H^2 = C_H(\mathbf{0})$ is the field variance. $\rho = 1$ means two points have identical values; $\rho \approx 0$ means they are nearly independent.

The rate at which $\rho$ decays is controlled by the \textbf{correlation length} $L_c > 0$. Two common models for concrete \cite{sudret2000, fenton2014}:

\begin{equation}
    \rho_H(\tau) =
    \exp\!\left(-\frac{\tau^2}{L_c^2}\right)
    \;\text{(squared-exponential)},
    \qquad
    \rho_H(\tau) =
    \exp\!\left(-\frac{|\tau|}{L_c}\right)
    \;\text{(exponential)}.
    \label{eq:corr-models}
\end{equation}

Physically, $L_c$ represents the distance beyond which material properties become essentially independent. A small $L_c$ means the material varies rapidly in space (highly heterogeneous); a large $L_c$ means properties change slowly (nearly uniform).

\paragraph{Remark:} In this work, $L_c$ is not treated as a parameter to be varied. A single representative value is adopted from experimental data \cite{delarrard2010} to account for spatial heterogeneity in the Monte Carlo simulations. The objective is not to study the effect of $L_c$ itself, but to compare the structural response \emph{with} and \emph{without} spatial variability.

\subsection{Random field simulation} \label{sec:rf-methods}
To incorporate a random field into a finite element model, one must generate realisations that respect the prescribed covariance structure. While several approaches exist \cite{fenton2014}, methods like Cholesky decomposition or Karhunen--Loève (KL) expansion scale poorly and become computationally prohibitive for large meshes. Conversely, efficient alternatives like FFT-based methods or Local Average Subdivision (LAS) are typically restricted to regular rectangular grids.

To overcome these limitations, this work adopts the \textbf{Turning Bands Method (TBM)} \cite{matheron1973, fenton2014}. The TBM simplifies the generation process by reducing the $d$-dimensional spatial simulation to a set of independent one-dimensional stationary processes along random lines (bands) passing through the domain. Because it requires neither a global matrix factorisation nor an eigenvalue solve, the TBM is computationally highly efficient and operates directly on any complex mesh geometry.


\subsection{Turning Bands Method (\texttt{ALEA})} \label{sec:tbm}
\paragraph{Principle}

The Turning Bands Method (TBM), introduced by \cite{matheron1973}, simplifies the generation of a 2-D or 3-D random field by reducing it to a set of 1-D problems.

The idea: draw $L$ lines through the domain in different directions. Along each line, simulate an independent 1-D random process. To obtain the field value at any point $\mathbf{x}$, project it onto each line and average the results:

\begin{equation}
    Z(\mathbf{x}) \;=\; \frac{1}{\sqrt{L}}
    \sum_{i=1}^{L}
    Z_i\!\left(\mathbf{x} \cdot \mathbf{u}_i\right),
    \label{eq:tbm}
\end{equation}

where $\mathbf{u}_i$ is the direction of the $i$-th line and $Z_i$ is the 1-D process along it. The 1-D processes are constructed so that the resulting field reproduces the prescribed covariance structure (\cref{sec:correlation}). In practice, $L = 16$ lines are sufficient for isotropic fields in two dimensions \cite{fenton2014}.

\paragraph{Implementation in Cast3M}

The TBM is natively implemented in the Cast3M \texttt{ALEA} operator, which takes as input the Gaussian parameters $(\mu_{\ln}, \sigma_{\ln})$ and the correlation length $L_c$, and returns the simulated field directly as nodal values on the mesh. The lognormal transformation and assignment to \texttt{MATE} follow the procedure described in \cref{sec:lognormal-transformation}.

\subsection{Transformation of Lognormal to Gaussian distribution} \label{sec:lognormal-transformation}

In practical finite element implementations using the Cast3M software, material parameters such as Young's modulus and permeability must remain strictly positive. However, the built-in random field operator \texttt{ALEA} generates Gaussian (normal) distributions, necessitating a mathematical transformation \cite{sudret2000}.

A \textbf{lognormal random field} $H(\mathbf{x})$ is therefore adopted: by definition, $\ln H(\mathbf{x})$ is a Gaussian random field, which guarantees $H(\mathbf{x}) > 0$ at every point \cite{delarrard2010}. The moments of $H(\mathbf{x})$ are related to those of the underlying Gaussian field $G(\mathbf{x}) \sim \mathcal{N}(\mu_{\ln},\,\sigma_{\ln}^2)$ by \cite{sudret2000}:

\begin{equation}
    \begin{cases}
        \mu_X \;=\; \exp\!\left(\mu_{\ln} + \dfrac{\sigma_{\ln}^2}{2}\right) \\[8pt]
        \sigma_X^2 \;=\; \left[\exp(\sigma_{\ln}^2) - 1\right]
                         \exp\!\left(2\mu_{\ln} + \sigma_{\ln}^2\right)
    \end{cases}
    \label{eq:lognormal-system}
\end{equation}

where $\mu_X$ and $\sigma_X$ are the physical mean and standard deviation of the material property. Inverting \cref{eq:lognormal-system} yields the Gaussian parameters as functions of the physical moments:

\begin{equation}
    \sigma_{\ln}^2 = \ln\!\left(1 + CoV_X^2\right),
    \qquad
    \mu_{\ln} = \ln(\mu_X) - \tfrac{1}{2}\,\sigma_{\ln}^2,
    \label{eq:lognormal-params}
\end{equation}

where $CoV_X = \sigma_X/\mu_X$ is the coefficient of variation.

The standard computational procedure in Cast3M \texttt{ALEA} then consists of two steps:

\begin{enumerate}
    \item Generate a Gaussian random field
          $G(\mathbf{x}) \sim \mathcal{N}(\mu_{\ln},\,\sigma_{\ln}^2)$ using the Turning Band Method (see \cref{sec:tbm}), with parameters computed from \cref{eq:lognormal-params}.

    \item Transform pointwise to obtain the lognormal field:
          \begin{equation}
              H(\mathbf{x}) \;=\;
              \exp\!\bigl(G(\mathbf{x})\bigr),
              \label{eq:lognormal-transform}
          \end{equation}
          which guarantees $H(\mathbf{x}) > 0$ and preserves the correlation structure characterised by $L_c$ (see \cref{sec:correlation}).
\end{enumerate}

This pointwise transformation is exact: the autocorrelation structure of the underlying Gaussian field is preserved in the lognormal field, so the spatial variability characterised by the correlation length $L_c$ remains physically meaningful after transformation \cite{delarrard2010}.

\section{Monte Carlo simulation and LHS framework} \label{sec:mc-lhs}
\subsection{Monte Carlo Principle}
Monte Carlo simulation (MCS) estimates the effect of uncertain inputs by running the deterministic model multiple times, each time with a different random sample drawn from the input distributions \cite{fenton2014}. However, standard MCS converges only as $1/\sqrt{N}$, necessitating a vast number of computationally expensive finite element runs.

\subsection{Latin Hypercube Sampling (LHS)} \label{sec:lhs-theory}
To reduce the computational cost and accelerate statistical convergence, \textbf{Latin Hypercube Sampling (LHS)} is utilized as a variance-reduced alternative. In LHS, the CDF of each random variable is divided into $N$ equiprobable strata (intervals of equal probability). Exactly one sample is drawn from each stratum by inverting the CDF, guaranteeing that the input space is covered evenly across the entire range of the probability distributions. This ensures the output statistics stabilize with markedly fewer evaluations compared to standard Monte Carlo.

\subsection{Convergence}

The probability of failure is estimated as the fraction of realisations where the limit state is violated:

\begin{equation}
    \hat{P}_f = \frac{N_f}{n},
    \label{eq:mc-pf}
\end{equation}

where the hat denotes that $\hat{P}_f$ is an estimate of the true $P_f$, which improves as $n$ increases, and $N_f$ is the number of failures out of $n$ total runs. The more runs, the more accurate the estimate---but the error only decreases as $1/\sqrt{n}$. To halve the error, one must quadruple the number of runs.

At $95\%$ confidence, the required number of runs is \cite{fenton2014}:

\begin{equation}
    n \;\geq\; \frac{4\,p_f(1-p_f)}{e^2},
    \label{eq:mc-nsample}
\end{equation}

where $e$ is the acceptable absolute error. For a small failure probability ($p_f \approx 10^{-3}$) with $10\%$ relative accuracy, this gives $n \approx 400{,}000$ runs. When a single THM analysis takes several minutes, this cost becomes prohibitive---motivating the Gaussian Quadrature method (\cref{sec:quadrature}) as a faster alternative.


\section{Gaussian Quadrature framework} \label{sec:quadrature}

Unlike Monte Carlo, which runs thousands of simulations, Gaussian Quadrature evaluates the model at a small number of strategically chosen points and reconstructs the output statistics from a weighted sum \cite{sudret2003}. The method has three steps, corresponding to three Cast3M operators:

\begin{enumerate}
    \item \textbf{\texttt{QUADRATU}}: compute the evaluation points and weights for each random variable.
    \item \textbf{\texttt{PARASTAT}}: run the model at each point combination and compute the first four statistical moments (mean, standard deviation, skewness, kurtosis).
    \item \textbf{\texttt{PROBDENS}}: reconstruct the PDF from these moments and estimate $P_f$ and $\beta$.
\end{enumerate}

\subsection{Step 1: Quadrature points and weights (\texttt{QUADRATU})} \label{sec:quad-points}

Computing any statistical moment of the output $G = g(\mathbf{X})$ requires evaluating an integral:

\begin{equation}
    I \;=\; \int_{\mathcal{D}} g(x)\, W(x)\,\mathrm{d}x,
    \label{eq:quad-integral}
\end{equation}

where $W(x) = f_X(x)$ is the PDF of the input variable. Gaussian quadrature approximates this by a weighted sum over $K$ evaluation points:

\begin{equation}
    I \;\approx\; \sum_{k=1}^{K} w_k\, g(x_k),
    \label{eq:quad-1d}
\end{equation}

where $\{x_k\}$ are the \emph{abscissas} (evaluation points) and $\{w_k\}$ the corresponding \emph{weights}. The optimal abscissas are the roots of orthogonal polynomials associated with $W$, and a $K$-point scheme integrates exactly any polynomial of degree up to $2K - 1$ \cite{sudret2003}.

\texttt{QUADRATU} automatically computes $\{x_k, w_k\}$ for any prescribed distribution (Normal, Lognormal, Uniform, etc.). The user provides only the distribution type and its parameters.

\subsection{Step 2: Statistical moments (\texttt{PARASTAT})} \label{sec:quad-moments}

For $N$ independent random variables with $K$ points each, the quadrature is applied to each variable separately and combined as a tensor product \cite{sudret2003}:

\begin{equation}
    \mu_q^{(0)} \;\approx\;
    \sum_{k_1=1}^{K_1} \cdots \sum_{k_N=1}^{K_N}
    w_{k_1} \cdots w_{k_N}\,
    \bigl[g(x_{k_1}, \ldots, x_{k_N})\bigr]^q,
    \label{eq:moment-quadrature}
\end{equation}

requiring $K^N$ model evaluations in total. For example, with $N = 2$ variables and $K = 3$ points each: $3^2 = 9$ runs.

\texttt{PARASTAT} collects these evaluations and returns the first four moments: mean $\mu$, standard deviation $\sigma$, skewness $\delta_3 = \mu_3/\sigma^3$, and kurtosis $\delta_4 = \mu_4/\sigma^4$.

\subsection{Step 3: PDF reconstruction and reliability (\texttt{PROBDENS})} \label{sec:quad-pdf}
Once the four moments are known, \texttt{PROBDENS} reconstructs the PDF using \emph{three independent methods simultaneously} \cite{sudret2003}:
\begin{itemize}
    \item \textbf{Winterstein approximation:} expresses $G$ as a Hermite polynomial transformation of a standard normal variable. Simplest but least accurate of the three methods.
    \item \textbf{Pearson's method:} selects the best-fit distribution family (normal, beta, gamma, etc.) from the moment values. Most accurate, but does not always provide a result.
    \item \textbf{Exponential of polynomials ($n=4$):} models the PDF as $\tilde{f}_G = c\,\exp(Q)$ where $Q$ is a polynomial fitted to the moments. Always provides accurate results.
\end{itemize}

The reconstructed Cumulative Distribution Function (CDF) serves as the fundamental basis for calculating the probability of failure $P_f$. Each method independently yields a probability of failure, obtained directly from the reconstructed CDF $\tilde{F}_G$ as:
\begin{equation}
    P_f \;=\; \tilde{F}_G(0),
    \label{eq:pf-quad}
\end{equation}
from which the reliability index follows through the general definition (\cref{eq:beta-generalized}). When the three reconstructions are consistent, the result can be accepted with confidence; the method is accurate for $\beta \lesssim 5$ \cite{sudret2003}.

\subsection{Computational cost} \label{sec:quad-cost}

The total number of model evaluations is $K^N$, which grows exponentially with $N$. For the present study with $N = 2$ and $K = 3$, this gives only 9 runs---compared to thousands for Monte Carlo. However, for large $N$ (e.g.\ $3^{10} = 59{,}049$ runs), the cost becomes comparable to Monte Carlo. This \emph{curse of dimensionality} limits the method to problems with a small number of random variables ($N \lesssim 10$--$15$) \cite{sudret2003}.

\section{Sensitivity indicators and ranking methods} \label{sec:sensitivity-indicators}
\subsection{Variance decomposition} \label{sec:var-decomp}

For a model $G = g(X_1, \ldots, X_N)$ with $N$ independent random inputs, Hoeffding's decomposition expresses $G$ as a sum of functions of increasing dimension:

\begin{equation}
    G \;=\; g_0
    \;+\; \sum_{i=1}^{N} g_i(X_i)
    \;+\; \sum_{i<j} g_{ij}(X_i, X_j)
    \;+\; \cdots
    \;+\; g_{1,\ldots,N}(X_1,\ldots,X_N),
    \label{eq:hoeffding}
\end{equation}

where $g_0 = \mathbb{E}[G]$ is a constant and each term has zero mean. Taking the variance of both sides yields:

\begin{equation}
    \mathrm{Var}(G) \;=\;
    \sum_{i=1}^{N} V_i
    \;+\; \sum_{i<j} V_{ij}
    \;+\; \cdots
    \;+\; V_{1,\ldots,N},
    \label{eq:var-decomp}
\end{equation}

where $V_i = \mathrm{Var}[g_i(X_i)]$ is the first-order variance (effect of $X_i$ alone) and $V_{ij} = \mathrm{Var}[g_{ij}(X_i,X_j)]$ is the second-order variance (interaction between $X_i$ and $X_j$).

\subsection{Variance-based sensitivity: Sobol indices} \label{sec:sobol}

Dividing each term in \cref{eq:var-decomp} by the total variance defines the \emph{Sobol indices}:

\begin{equation}
    S_i \;=\; \frac{V_i}{\mathrm{Var}(G)},
    \qquad
    S_{ij} \;=\; \frac{V_{ij}}{\mathrm{Var}(G)},
    \label{eq:sobol-indices}
\end{equation}

where:

\begin{itemize}
    \item $S_i$ is the \textbf{first-order Sobol index}: fraction of the output variance due to $X_i$ alone.
    \item $S_{ij}$ is the \textbf{second-order Sobol index}: fraction due to the interaction between $X_i$ and $X_j$, beyond their individual effects.
\end{itemize}

The \textbf{total-order Sobol index} captures the full effect of $X_i$, including all interactions:

\begin{equation}
    S_{T_i} \;=\; S_i
    \;+\; \sum_{j \neq i} S_{ij}
    \;+\; \cdots
    \label{eq:sobol-total}
\end{equation}

By construction, the first-order indices sum to at most one ($\sum_i S_i \leq 1$), and the total-order indices sum to at least one ($\sum_i S_{T_i} \geq 1$). When the model is additive (no interactions), $S_{ij} = 0$ for all pairs and $S_i = S_{T_i}$.

The first-order index $S_i$ can be expressed as a conditional variance:

\begin{equation}
    S_i \;=\;
    \frac{\mathrm{Var}_{X_i}\bigl[
          \mathbb{E}_{X_{\sim i}}(G \mid X_i)\bigr]}
         {\mathrm{Var}(G)},
    \label{eq:sobol-cond}
\end{equation}

where $X_{\sim i}$ denotes all variables except $X_i$. Intuitively, $\mathbb{E}(G \mid X_i)$ averages out all other variables: if this conditional mean varies a lot with $X_i$, then $X_i$ is influential.

A variable with a large $S_i$ (or $S_{T_i}$) is a \emph{dominant parameter}: it drives most of the output uncertainty and should be characterised with care. A variable with a small index could be treated as deterministic without significant loss of accuracy.

\section{Probability of failure and reliability index} \label{sec:pf-beta}
\subsection{Limit state function} \label{sec:limit-state}
In structural reliability, a \textbf{performance function} (or limit-state function) $G(\mathbf{X})$ explicitly separates the structural state into two domains:
\begin{equation}
    G(\mathbf{X}) > 0 \Rightarrow \text{safety domain } D_s, \qquad G(\mathbf{X}) \leq 0 \Rightarrow \text{failure domain } D_f.
    \label{eq:limit-state}
\end{equation}
For a typical structure evaluated against a specific load $S$ and a critical resistance $R$, the function reads $G(R,S) = R - S$ \cite{baroth2011}.

\subsection{Probability of failure and reliability index $\beta$}
The \textbf{probability of failure} is the integral of the joint probability density function over the failure domain \cite{baroth2011}:
\begin{equation}
    P_f = \mathrm{Prob}\bigl[G(\mathbf{X}) \leq 0\bigr] = \int_{G(\mathbf{X}) \leq 0} f_{\mathbf{X}}(\mathbf{x}) \,\mathrm{d}\mathbf{x}.
    \label{eq:pf}
\end{equation}

However, when dealing with complex non-linear finite element models, this multidimensional integral cannot be solved analytically and must be estimated through numerical approximation methods such as Monte Carlo simulation or Kernel Density Estimation (KDE).

In modern engineering codes, $P_f$ is frequently mapped to the \textbf{generalized reliability index} $\beta$ via the inverse of the standard normal cumulative distribution function $\Phi^{-1}$:
\begin{equation}
    \beta = -\Phi^{-1}(P_f).
    \label{eq:beta-generalized}
\end{equation}

\section{Reference parameters for the simplified poro-mechanical model} \label{sec:simplified-model-inputs}
To evaluate the effect of spatial heterogeneity in \cref{chapter5}, a simplified isothermal poro-mechanical model is implemented in \textsc{Cast3M}. This section details the specific constitutive laws and material parameters used for that test case.

\subsection{Constitutive relationships}
The hygral behaviour is governed by the van Genuchten model. The degree of saturation $S_l$ and the liquid relative permeability $k_{rl}$ are functions of the capillary pressure $P_c$:
\begin{equation}
    S_l(P_c) = \left[ 1 + \left(\frac{P_c}{a_1}\right)^{\frac{b_1}{b_1-1}} \right]^{-\frac{1}{b_1}}
\end{equation}
\begin{equation}
    k_{rl}(P_c) = \sqrt{S_l} \left[ 1 - \left(1 - S_l^{b_1}\right)^{\frac{1}{b_1}} \right]^2
\end{equation}
where $a_1$ and $b_1$ are material parameters. The effective hydraulic conductivity $K_{\mathrm{eff}}$ entering the flow matrix is:
\begin{equation}
    K_{\mathrm{eff}}(P_c) = \frac{k_{rl}(P_c)\,\rho_l\,K_0}{\mu_l}
\end{equation}

\subsection{Geometry, mesh, and material parameters}
The modelled specimen is a $10 \times 10 \times 10$~cm cube. By exploiting symmetry, only 1/8 of the domain is meshed using 3D 8-node hexahedral elements (\texttt{CUB8}). The mesh is refined near the exposed surfaces, transitioning from a coarse element size of $10$~mm down to $2$~mm. The transient simulation spans $120$ days with a time step of $1$ day.

The deterministic reference parameters fed into the model are summarised in \cref{tab:simplified-params}. For the stochastic analyses (C1, C2, C3), the intrinsic permeability $K_0$ is treated as a random variable or a spatial random field centered around the mean value of $1.0\times 10^{-20}~\mathrm{m^2}$.

\begin{table}[htbp]
\centering
\caption{Reference parameters for the simplified isothermal poro-mechanical model.}
\label{tab:simplified-params}
\begin{tabular}{lllc}
\toprule
Parameter & Symbol & Value & Unit \\
\midrule
\multicolumn{4}{l}{\emph{Hygral properties}} \\
Liquid water density & $\rho_l$ & $1000$ & $\mathrm{kg/m^3}$ \\
Dynamic viscosity & $\mu_l$ & $1.0 \times 10^{-3}$ & $\mathrm{Pa \cdot s}$ \\
Porosity & $\phi$ & $0.20$ & $-$ \\
Mean intrinsic permeability & $K_0$ & $1.0 \times 10^{-20}$ & $\mathrm{m^2}$ \\
van Genuchten parameter & $a_1$ & $20$ & $\mathrm{MPa}$ \\
van Genuchten parameter & $b_1$ & $2.0$ & $-$ \\
\addlinespace
\multicolumn{4}{l}{\emph{Mechanical properties}} \\
Young's modulus & $E$ & $30$ & $\mathrm{GPa}$ \\
Poisson's ratio & $\nu$ & $0.20$ & $-$ \\
Biot coefficient & $\alpha$ & $0.60$ & $-$ \\
\addlinespace
\multicolumn{4}{l}{\emph{Initial and boundary conditions}} \\
Temperature & $T$ & $20$ & $^\circ\mathrm{C}$ \\
Initial relative humidity & $\mathrm{RH}_{\mathrm{ini}}$ & $0.90$ & $-$ \\
External relative humidity & $\mathrm{RH}_{\mathrm{ext}}$ & $0.50$ & $-$ \\
\bottomrule
\end{tabular}
\end{table}